\chapter{Einleitung} \label{ch:einleitung}

%\section{Problemstellung}
Bei der Suche nach einem Cookie-Rezept zeigt sich, dass bei vielen Rezepten \footnote{Für die Rezepte siehe Anhang \ref{ch:appendix_a} und deren Links in Anhang \ref{ch:appendix_b}.} 
die Mengenangaben im Verhältnis zur Mehlmenge sehr ähnlich sind (siehe Tabelle \ref{table:rezepte_relativ}).
Die Menge an Zucker und Butter unterscheiden sich allerdings stark.
Für 100\,g Mehl varriert die Buttermenge von 48,00\,g bis 89,29\,g,
die Menge des Zuckers von 67,20\,g bis 115,56\,g.
Bei dem breiten Wertebereich fragt sich, wie deren Menge sich auf den Cookie auswirkt.
Auch die Menge an Schokoladenstückchen (SchokoStü) variiert in den Rezepten werden in dieser Arbeit aber nicht weiter untersucht.
Die Werte für Natron, Eier und Salz sind relativ konstant.
Pareyt \etal \cite{pareyt2009role} zeigten bereits eine Abhängigkeit der Cookie-Struktur von Zucker- und Buttermenge.
Allerdings fehlt hier die Ermittlung der Faktorwechselwirkungen (FWW) für die Zuckermenge, Buttermenge und deren Zusammenspiel.

\begin{table}[h!]
\centering
\caption{Auf 100\,g Mehl normierte Zutatenmengen der ausgewerteten Rezepte}
\label{table:rezepte_relativ}
\footnotesize
\setlength{\tabcolsep}{4pt}
\begin{tabular}{|l|r|r|r|r|r|r|}
\hline
Rezept & Butter (g) & Zucker (g) & Eier (St.) & Natron (g) & Salz (g) & SchokoSt (g) \\
\hline
R1     & 89,29      & 92,86      & 0,71       & 1,07       & 0,89     & 107,14    \\
R2     & 89,29      & 85,00      & 0,71       & 1,07       & 0,89     & 71,43     \\
R3     & 68,00      & 67,20      & 0,60       & 0,60       & 0,68     & 80,00     \\
R4     & 60,71      & 92,14      & 0,54       & 1,07       & 0,18     & 53,57     \\
R5     & 66,67      & 100,00     & 0,53       & 1,07       & 0,13     & 80,00     \\
R6     & 70,83      & 83,33      & 0,42       & 1,25       & 1,04     & 104,17    \\
R7     & 72,22      & 115,56     & 0,56       & 0,83       & 0,28     & 77,78     \\
R8     & 48,00      & 73,60      & 0,40       & 1,20       & 0,20     & 60,00     \\
\hline
\end{tabular}
\end{table}

% 1TL (gestr.) Salz entspricht 2,5 g
% 1TL Natron entspricht 3 g
% 1 Pck. Vanillezucker entspricht 8 g
% 1 TL Vanillezucker entspricht 4 g
% 1 Prise Salz entspricht 0,5 g
% Eier liegen meist bei 2 Stück
% Eier manchmal auch ein ganzen und ein Eigelb

% R1 : https://www.chefkoch.de/rezepte/2257651361176625/Subway-Cookies.html  
% R2 : https://www.chefkoch.de/rezepte/1378031243002711/American-Cookies-wie-bei-Subway.html  
% R3 : https://www.chefkoch.de/rezepte/470791140633601/Chewy-Chocolate-Chip-Cookies.html  
% R4 : https://www.chefkoch.de/rezepte/2601991408817771/American-Soft-Chocolate-Chip-Cookies.html  
% R5 : https://www.chefkoch.de/rezepte/614921161520963/World-s-best-Chocolate-Chip-Cookies.html  
% R6 : https://cookidoo.de/recipes/recipe/de-DE/r307943  
% R7 : https://cookidoo.de/recipes/recipe/de-DE/r47866    
% R8 : https://cookidoo.de/recipes/recipe/de-DE/r136546  



Beim Zubereiten von Cookies mit Vollkornanteilen erhöhte die ernährungsphysiologischen Eigenschaften deutlich, wie Unal und Ozkaya \cite{arslan2024nutritional} nachwiesen.
Dabei untersuchten sie den Vollkornmehltyp und - anteil, die Wechselwirkungen mit Zucker- und Buttermenge blieben dabei offen.
%Gleichzeitig verändert Vollkornmehl oft sensorische und technologische Eigenschaften, sodass Rezepturanpassungen notwendig werden können.
%Rosen \etal \cite{rosen2008gradual} berichten über die Schrittweise Erhöhung des Vollkornanteils in Brotprodukten.
%Sie erhöhten damit die Vollkornaufnahme bei Kindern und verbesserten die Akzeptanz gegenüber einer direkten 100\,\% Vollkorn Umstellung.  

%\section{Zielsetzung}\label{sec:zielsetzung}

Diese Arbeit schließt diese Lücken.
Sie ermittelt die Wechselwirkungen von Zucker, Butter und Vollkornmehl in Cookies.
Zusätzlich prüft sie die Eignung eines linearen Modells im betrachteten Faktorraum.
Dafür formuliert die Arbeit folgende Forschungsfrage (FF):

\begin{quote}
    Wie verändern Zucker, Butter und Vollkornmehlanteil gemeinsam die geometrischen Eigenschaften von Keksen, und treten dabei Wechselwirkungen auf?
\end{quote}


% Nutzung des 2³-Vollfaktorplan mit Zentralpunkten und Blöcken.
% Charakterisierung des Cookiedurchmesser in Abhängigkeit von der Butter- und Zuckermenge, sowie vom Vollkornmehlanteil.